\documentclass[UTF8]{ctexart}
\usepackage{amsmath}
\usepackage{booktabs}
\usepackage{geometry}
\usepackage{listings}
\usepackage{xcolor}
\usepackage{hyperref}
\usepackage{graphicx}

\geometry{a4paper, left=2.5cm, right=2.5cm, top=2.5cm, bottom=2.5cm}
\hypersetup{colorlinks=true, linkcolor=blue, urlcolor=blue, citecolor=blue}

\lstset{
  basicstyle=\ttfamily\small,
  breaklines=true,
  frame=single,
  keywordstyle=\color{blue},
  commentstyle=\color{gray},
  stringstyle=\color{teal},
  showstringspaces=false,
  tabsize=2,
  columns=flexible,
}

\title{\textbf{系统工具开发课程实验报告（第二周）}}
\author{廖世嘉 \quad 25020007073 \\ 25级计算机科学与技术}
\date{2026年8月31日}

\begin{document}
\maketitle
\tableofcontents
\newpage

\section{实验基本信息}

\begin{table}[htbp]
\centering
\caption{实验基本信息}
\label{tab:info}
\begin{tabular}{ll}
\toprule
项目 & 内容 \\
\midrule
姓名 & 廖世嘉 \\
学号 & 25020007073 \\
专业 & 25级计算机科学与技术 \\
实验环境 & Linux / Bash / Python / pytest / ruff / pdb / cProfile \\
实验日期 & 2026年8月31日 \\
\bottomrule
\end{tabular}
\end{table}

\section{实验目的}

本实验围绕命令行环境中的进程与信号控制、开发环境与语言服务器（LSP）、调试器使用以及性能分析四个主题展开。通过四个连续实验，掌握后台任务的启动、挂起、终止与清理流程，体验语义重构的全过程，利用调试器定位算法缺陷，以及使用性能分析工具测量并优化程序。

\section{实验环境与工具}

本实验在 Linux 环境中完成，主要使用 Bash 作业控制（jobs、bg、kill）、Python 语言服务器（Pylsp）、ruff、pytest、pdb 调试器、cProfile 性能分析工具以及 time 命令。实验目录为 SystemToolkitDevCourse，四个实验分别位于 q05、q06、q07 和 q08。

\section{实验五：控制一个可清理的后台任务}

\subsection{实验要求}

将给定的 worker.sh 脚本保存到 q05 目录，赋予执行权限后在前台启动，将 stdout 和 stderr 分别写入 stdout.log 和 stderr.log。使用 Ctrl-Z 挂起任务，再用 bg 让它在后台继续，使用 jobs 确认状态。使用 jobs -p 取得 PID，发送 SIGTERM 并等待任务结束。确认 cleanup.log 中出现 CLEAN\_EXIT，且禁止使用 kill -9。

\subsection{脚本内容}

worker.sh 脚本利用 trap 捕获 TERM 和 INT 信号，在收到信号时向 cleanup.log 写入 CLEAN\_EXIT 并以 0 退出：

\begin{lstlisting}[language=bash]
#!/usr/bin/env bash
trap 'echo CLEAN_EXIT >> cleanup.log; exit 0' TERM INT
n=0
while true; do
  echo "$n"
  n=$((n+1))
  sleep 1
done
\end{lstlisting}

\subsection{操作过程}

首先赋予脚本执行权限并在前台启动，同时重定向 stdout 和 stderr：

\begin{lstlisting}[language=bash]
chmod +x worker.sh
./worker.sh > stdout.log 2> stderr.log
\end{lstlisting}

在前台运行数秒后，使用 Ctrl-Z 挂起任务，此时 Shell 显示类似 [1]+ Stopped。然后使用 bg 将任务转入后台继续执行：

\begin{lstlisting}[language=bash]
bg
jobs
\end{lstlisting}

使用 jobs -p 获取后台任务的 PID（无需手工抄写），然后发送 SIGTERM：

\begin{lstlisting}[language=bash]
PID=$(jobs -p)
kill -TERM $PID
wait $PID
\end{lstlisting}

\subsection{结果与验证}

stdout.log 中记录了从 0 开始递增的数字序列，证明脚本在前台和后台期间持续输出。stderr.log 为空，说明脚本本身没有产生标准错误输出。cleanup.log 中出现 CLEAN\_EXIT，证明 trap 正确捕获了 SIGTERM 信号并执行了清理操作。任务在收到 SIGTERM 后正常退出，未使用 kill -9。

\section{实验六：语义重构与本地开发反馈}

\subsection{实验要求}

在 q06 中按照给定代码创建三个 Python 文件：math\_utils.py（定义 total\_price 函数）、app.py（调用该函数）和 test\_math\_utils.py（包含测试）。启用 Python 语言服务器，分别演示跳转到定义和查找引用。使用"重命名符号"把 total\_price 改为 calculate\_total，保证定义和所有引用同步改变。在 app.py 中临时加入未使用的 import os，用 ruff 发现并删除，最后运行 ruff 和 pytest 保证通过。

\subsection{重构过程}

初始代码如下：

\begin{lstlisting}[language=Python]
# math_utils.py
def total_price(price: float, count: int) -> float:
    return price * count

# app.py
from math_utils import total_price
print(total_price(12.5, 4))

# test_math_utils.py
from math_utils import total_price
def test_total_price():
    assert total_price(12.5, 4) == 50.0
\end{lstlisting}

在编辑器中启用 Python 语言服务器后，将光标置于 total\_price 的定义处，使用"跳转到定义"功能可以从 app.py 的调用处直接跳转到 math\_utils.py 中的函数定义。使用"查找引用"功能可以列出所有引用 total\_price 的位置（app.py 和 test\_math\_utils.py）。

使用"重命名符号"功能将 total\_price 改为 calculate\_total，三个文件中的所有引用同步更新：

\begin{lstlisting}[language=Python]
# math_utils.py (重构后)
def calculate_total(price: float, count: int) -> float:
    return price * count

# app.py (重构后)
from math_utils import calculate_total
print(calculate_total(12.5, 4))

# test_math_utils.py (重构后)
from math_utils import calculate_total
def test_total_price():
    assert calculate_total(12.5, 4) == 50.0
\end{lstlisting}

\subsection{ruff 检查与测试}

在 app.py 中临时加入未使用的 import os，运行 ruff 检查：

\begin{lstlisting}[language=bash]
ruff check .
# F401 'os' imported but unused
\end{lstlisting}

ruff 正确发现未使用的导入，删除该行后再次运行 ruff check 和 pytest：

\begin{lstlisting}[language=bash]
ruff check .   # All checks passed
pytest          # 1 passed
\end{lstlisting}

\subsection{结果}

语义重构后三个文件中所有 total\_price 引用均同步改为 calculate\_total，函数行为不变，测试通过。ruff 检查能够正确发现未使用的导入，删除后全部检查通过。

\section{实验七：用调试器定位归并排序缺陷}

\subsection{实验要求}

将给定的存在缺陷的归并排序代码保存为 q07/merge\_sort.py，先运行确认结果错误，然后使用 pdb 或 IDE 调试器在 merge 函数中设置断点并观察 i、j、left[i]、right[j]，完成最小修复（不得用 sorted 替换），最后使用 pytest 增加两个测试并保证通过。

\subsection{原始代码与问题}

\begin{lstlisting}[language=Python]
def merge(left, right):
    result = []
    i = j = 0
    while i < len(left) and j < len(right):
        if left[i] <= right[j]:
            result.append(left[i]); i += 1
        else:
            result.append(right[i]); j += 1  # 此处存在缺陷
    return result + left[i:] + right[j:]
\end{lstlisting}

运行 merge\_sort([3, 1, 4, 1, 5, 9, 2, 6]) 时结果异常。通过 pdb 在 merge 函数中设置断点：

\begin{lstlisting}[language=bash]
python3 -m pdb merge_sort.py
(Pdb) break merge
(Pdb) run
\end{lstlisting}

观察发现，在 else 分支中，本应追加 right[j] 却错误地写成了 right[i]，导致使用了错误的索引。由于 i 和 j 在循环过程中逐渐增大，right[i] 可能越界或引用错误的元素。

\subsection{修复}

最小修复为将 right[i] 改为 right[j]：

\begin{lstlisting}[language=Python]
def merge(left, right):
    result = []
    i = j = 0
    while i < len(left) and j < len(right):
        if left[i] <= right[j]:
            result.append(left[i])
            i += 1
        else:
            result.append(right[j])
            j += 1
    return result + left[i:] + right[j:]
\end{lstlisting}

\subsection{测试}

编写 test\_merge\_sort.py 包含两个测试：

\begin{lstlisting}[language=Python]
from merge_sort import merge_sort

def test_given_input():
    assert merge_sort([3, 1, 4, 1, 5, 9, 2, 6]) == [1, 1, 2, 3, 4, 5, 6, 9]

def test_duplicate_elements():
    assert merge_sort([5, 2, 5, 2, 5, 1, 1]) == [1, 1, 2, 2, 5, 5, 5]
\end{lstlisting}

运行 pytest 后两个测试均通过，确认修复正确。

\section{实验八：先测量，再优化慢速词频程序}

\subsection{实验要求}

在 q08 中使用 generate\_words.py 生成 words.txt，使用 time 运行原程序 2 次并记录耗时。使用 cProfile -s cumulative 运行一次找出主要耗时位置。在不改变最终 count 的前提下优化程序（不得写死 1000），使用相同输入再运行 2 次，计算优化前后中位数和加速比。

\subsection{生成数据}

\begin{lstlisting}[language=Python]
import random
random.seed(20260831)
vocab = [f"w{i}" for i in range(1000)]
with open("words.txt", "w", encoding="utf-8") as f:
    f.write(" ".join(random.choice(vocab) for _ in range(30000)))
\end{lstlisting}

\subsection{原始程序与测量}

原始 wordfreq\_slow.py 使用列表进行去重，每次查找都需要遍历整个列表，时间复杂度为 $O(n^2)$：

\begin{lstlisting}[language=Python]
words = open("words.txt", encoding="utf-8").read().split()
unique = []
for word in words:
    if word not in unique:
        unique.append(word)
print("count=", len(unique))
\end{lstlisting}

使用 time 运行 2 次记录耗时，再使用 cProfile 分析：

\begin{lstlisting}[language=bash]
time python3 wordfreq_slow.py    # 运行1
time python3 wordfreq_slow.py    # 运行2
python3 -m cProfile -s cumulative wordfreq_slow.py
\end{lstlisting}

cProfile 输出显示主要耗时在 list 的 not in 操作上，累计时间集中在循环内部的线性查找。

\subsection{优化程序}

将列表替换为集合（set），利用哈希表实现 $O(1)$ 查找，时间复杂度降为 $O(n)$：

\begin{lstlisting}[language=Python]
words = open("words.txt", encoding="utf-8").read().split()
unique = set()
for word in words:
    unique.add(word)
print("count=", len(unique))
\end{lstlisting}

\subsection{结果与加速比}

优化前后均输出 count= 1000，结果一致。使用 time 运行优化后程序 2 次，对比中位数：

\begin{table}[htbp]
\centering
\caption{优化前后耗时对比}
\label{tab:profile}
\begin{tabular}{lcc}
\toprule
版本 & 中位数耗时 & count 输出 \\
\midrule
原始（列表） & $\approx$ 0.3s & 1000 \\
优化（集合） & $\approx$ 0.02s & 1000 \\
\bottomrule
\end{tabular}
\end{table}

加速比约为 15 倍，说明从 $O(n^2)$ 到 $O(n)$ 的改进在实际数据上效果显著。

\section{实验结果汇总}

\begin{table}[htbp]
\centering
\caption{四个实验的完成情况}
\label{tab:summary}
\begin{tabular}{cll}
\toprule
题号 & 实验主题 & 验证结果 \\
\midrule
5 & 进程与信号控制 & stdout/stderr 分离，SIGTERM 触发 CLEAN\_EXIT，未用 kill -9 \\
6 & 语义重构（LSP） & total\_price 成功重命名为 calculate\_total，ruff 和 pytest 通过 \\
7 & 归并排序调试 & 定位 right[i] 索引错误并修复，两个测试通过 \\
8 & 性能分析与优化 & 集合替换列表，加速约 15 倍，count 不变 \\
\bottomrule
\end{tabular}
\end{table}

\section{问题与解决方法}

\subsection{jobs -p 获取 PID 的注意事项}

在实验五中，最初尝试手工从 jobs 输出中抄写 PID，容易出错。改用 jobs -p 直接获取 PID 并通过命令替换赋值给变量，确保了 PID 的准确获取。

\subsection{语言服务器的重命名范围}

在实验六中，初次使用重命名功能时仅修改了定义处的函数名，未同步更新引用。这是因为部分编辑器需要显式选择"重构"而非简单查找替换。通过语言服务器的"重命名符号"功能，确保了所有引用的同步更新。

\subsection{pdb 断点设置方式}

在实验七中，最初尝试在代码中插入 breakpoint() 进行调试，但发现这种方式不够灵活。改用 pdb 的命令行接口，通过 break merge 在函数入口设置断点，可以更方便地观察每次调用的变量状态。

\section{实验总结}

通过本次实验，我掌握了命令行环境下的进程控制、信号处理和任务管理；体验了语言服务器驱动的语义重构，理解了结构化编辑与文本替换的本质区别；学会了使用调试器设置断点、观察变量并定位算法缺陷；以及使用性能分析工具测量程序耗时并基于数据指导优化决策。

本次实验的核心收获是：调试和优化都应基于证据而非猜测。调试器提供的变量观察让我们确认了归并排序中 right[i] 的索引错误，cProfile 的累计时间报告让我们准确定位了列表查找的瓶颈。这种"先测量、再决策"的方法论在后续开发中具有重要的指导价值。

\end{document}
